\chapter{Phương pháp nghiên cứu\label{chapter:methodology}}
\chaptermark{Phương pháp}

Chương trước nêu khoảng trống. Chương này nêu cách nghiên cứu đóng nó.

\section{Phân hệ Quản lý Ngữ cảnh UE}\label{sec:uecm}
Phần được nghiên cứu là một dịch vụ Quản lý Ngữ cảnh UE. Nó được mô hình hoá theo một dịch vụ
gNB-CU-CP thực của Nokia, sở hữu vòng đời ngữ cảnh UE, thực thi thủ tục
RRC~\cite{3gpp38331}, và di động cùng chuyển giao. Đó là các chức năng mà thiết kế 3GPP giao
cho phần mặt điều khiển của một gNB tách (Mục~\ref{sec:ranarch}). Hình~\ref{fig:fsm} cho thấy
vòng đời ngữ cảnh UE mà cả hai bản mã hoá cùng một cách. Khối lượng hỗn hợp của
Mục~\ref{sec:workload} đưa mỗi UE đi qua máy trạng thái đó.

\begin{figure}[ht]
\centering
\begin{tikzpicture}[
    ->, >=Stealth, node distance=3.2cm, font=\small,
    every state/.style={draw, minimum size=1.5cm, fill=blue!6, thick},
    every edge/.style={draw, thick}
  ]
  \node[state, initial, initial text=\scriptsize tạo] (idle) {Idle};
  \node[state, right=of idle] (conn) {Connected};
  \node[state, right=of conn] (reconf) {Reconf.};
  \node[state, below=1.9cm of conn, fill=gray!10] (rel) {Released};

  \path (idle)  edge node[above,font=\scriptsize]{setup} (conn)
        (conn)  edge[bend left=18] node[above,font=\scriptsize]{begin} (reconf)
        (reconf) edge[bend left=18] node[below,font=\scriptsize]{complete / abort} (conn)
        (conn)  edge node[right,font=\scriptsize,pos=0.45]{release} (rel)
        (idle)  edge[bend right=32] node[below,font=\scriptsize,pos=0.6]{release} (rel);
  \path (conn) edge[loop above] node[font=\scriptsize]{signal / add DRB} (conn);
\end{tikzpicture}
\caption[Máy trạng thái vòng đời ngữ cảnh UE]{Máy trạng thái vòng đời ngữ cảnh UE.
Các thao tác không hợp lệ với trạng thái hiện tại trả về lỗi thay vì đổi
trạng thái.\label{fig:fsm}}
\end{figure}

UECM là công việc tầng 3 RAN cốt lõi. Nó có những đặc tính làm mã mặt điều khiển đang chạy trở
nên khó. Mỗi UE có trạng thái thay đổi. Vòng đời thao tác chặt. Nhiều UE chạy cùng lúc. Mục
tiêu trễ là những mục tiêu hệ thống 5G đặt cho thủ tục mặt điều
khiển~\cite{3gpp38913,3gpp22261}. Phát hiện về cách một ngôn ngữ xử lý những đặc tính đó liên
quan nhất tới các dịch vụ mặt điều khiển nhiều luồng, có trạng thái khác cùng loại. Phần này
cũng đủ nhỏ để đặc tả trọn vẹn và xây hai lần. Đó là điều làm một phép so sánh có kiểm thử
tương đương trở nên khả thi.

Giao diện nhỏ và được xây quanh vòng đời. Nó có thể tạo và giải phóng một ngữ cảnh. Nó có thể
thiết lập một kết nối. Nó có thể xử lý một thông điệp báo hiệu trong phiên. Nó có thể thêm một
data radio bearer. Nó có thể bắt đầu, hoàn tất, hoặc huỷ một tái cấu hình. Nó có thể tra trạng
thái của một ngữ cảnh dưới tải. Mỗi thao tác được định nghĩa cho mỗi trạng thái của
Hình~\ref{fig:fsm}. Một thao tác bất hợp pháp trả về lỗi có kiểu và để ngữ cảnh không đổi. Tập
đóng đó cho phép cùng các kiểm thử và cùng khối lượng đưa cả hai bản đi theo cùng các đường.

RQ2 cần một danh mục, không phải một danh sách mối nguy lỏng. Một danh mục được xây từ tập thao
tác này: $\FaultClasses$ lớp trong hai nhóm bảy. Bảy lớp đến từ bên trong vùng mà Rust an toàn
tuyên bố cai quản. Bảy lớp được chọn vì Rust an toàn được kỳ vọng không ngăn chúng. Mỗi lớp có
kỳ vọng được ghi trước khi các chương trình chạy. Mục~\ref{sec:safetyresults} báo cáo mỗi lớp
đã làm gì.

Đuôi là chỉ số vì hai đặc tính của phần này. Một hành động người dùng nhìn thấy, thiết lập cuộc
gọi hoặc chuyển giao, thành nhiều thao tác mặt điều khiển phụ thuộc nhau. Thành viên chậm nhất
chặn hoàn tất. Một site bận đưa hàng nghìn thủ tục như vậy mỗi giây. Vì vậy một điểm đuôi xa
được chạm đủ thường để có ý nghĩa. Dạng thất bại là lý do thứ hai. Một dịch vụ thiết kế này
theo tải đưa vào với độ trễ phẳng cho đến khi một luồng xử lý bão hoà. Rồi độ trễ nhảy mạnh
trong một bước tải. Đó là lý do nghiên cứu này quét qua bão hoà.

\section{Hướng tiếp cận nghiên cứu}\label{sec:approach}
Nghiên cứu chạy thí nghiệm. Nó cũng xây thứ nó đo: một cặp bản khớp nhau, cộng khung đo đưa
chúng chạy. Cả hai được xây từ một đặc tả cố định. Bằng chứng đến từ ba phương pháp tách biệt.
Một quét tải trả lời RQ1 và RQ3. Tiêm lỗi có kiểm soát trả lời RQ2. Một nguyên mẫu di trú FFI
trả lời RQ4. Một điểm yếu ở một trong số chúng, chẳng hạn máy chủ có thể dịch chuyển đuôi bao
nhiêu, không lan sang phần còn lại của các tuyên bố.

Mỗi câu hỏi hỏi một kết quả phụ thuộc vào một biến, ngôn ngữ. Ngôn ngữ phải là khoảng cách có
kế hoạch duy nhất giữa hai đối tượng thử. Điều đó loại trừ cách so sánh ngôn ngữ phổ biến nhất.
Ở dạng đó hai chương trình viết độc lập được đo, và khoảng cách được lấy là ngôn ngữ. Khoảng
cách đó cũng mang thiết kế, thư viện, và kỹ năng của tác giả. Xây cả hai phía từ một đặc tả loại
nhiễu loạn đó. Ở nơi câu hỏi trả lời được bằng số, nó dùng số. Ở nơi không thể, câu trả lời vẫn
theo một kế hoạch cố định, bằng lời. Đó là cách nghiên cứu này xử lý loại bảo đảm trong RQ2.
RQ4 hỏi về chi phí thiết kế của giao diện lập trình ứng dụng (API).

Công trình trước gần nhất là năm luận văn do nhà cung cấp lưu trữ của
Mục~\ref{sec:gap}. Chúng không xuất hiện trong cơ sở dữ liệu trích dẫn nào. Chúng phải được tìm
bằng cách tra kho lưu trữ của từng trường bằng tay. Cùng các tìm kiếm đó không tìm thấy luận
văn Phần Lan nào về Rust trong bối cảnh truy nhập vô tuyến. Chúng không tìm thấy bài được lập
chỉ mục nào nhắc Rust từ nhà cung cấp chủ trì công trình này. Tìm kiếm không trả về nghiên cứu
nào đo đuôi của một bản Rust và một bản C++ của cùng một phần mặt điều khiển nhiều luồng, có
trạng thái, dưới tải quá lớn.

Bảng~\ref{tab:design} ánh xạ mỗi câu hỏi nghiên cứu tới phương pháp, hiện vật, dữ liệu và chỉ
số trả lời nó. RQ1 và RQ3 dùng chung H1, diễn đạt để phủ một hoà tốc độ theo cả nghĩa trễ và
nghĩa tài nguyên. Cách diễn đạt đó hoá ra là điểm yếu chính của giả thuyết. RQ2 không phụ thuộc
máy chủ. RQ4 phụ thuộc theo một trục. Thời gian biên dịch là một phép đo thời gian. Vì vậy nó
được đọc như tỉ số giữa hai chuỗi công cụ đo luân phiên trên một máy.

\begin{table}[ht]
\centering
\small
\caption[Cách mỗi câu hỏi nghiên cứu được trả lời]{Cách mỗi câu hỏi nghiên cứu được trả lời.
RQ1 và RQ3 mang hai phương pháp chứ không phải một: quét đo vòng khứ hồi mà một triển khai sẽ
trải, và tách lõi khỏi bàn giao đo phần mà ngôn ngữ chịu trách nhiệm.
\label{tab:design}}
\renewcommand{\arraystretch}{1.2}
\setlength{\tabcolsep}{4pt}
\begin{tabular}{@{}p{0.07\linewidth} p{0.27\linewidth} p{0.27\linewidth} p{0.27\linewidth}@{}}
\hline\hline
RQ & Phương pháp & Hiện vật / dữ liệu & Chỉ số chính \\
\hline
RQ1 & Quét tải tới quá tải, lặp đan xen & Khung đo, histogram độ trễ &
trung vị, p99/p99.9/p99.99, và cực đại so với tải đưa vào \\
RQ3 & Cùng quét, thu tài nguyên & Mẫu \texttt{getrusage} mỗi lần chạy &
kích thước tập trú (RSS) đỉnh, CPU, thông lượng \\
RQ1, RQ3 & Gọi lõi trực tiếp, rồi qua dịch vụ &
Cặp microbenchmark khớp &
thao tác đối với bàn giao giữa luồng \\
RQ2 & Tiêm lỗi có kiểm soát & Danh mục lỗi, ba bộ dò lỗi C++,
\texttt{rustc}, Loom/Shuttle, Miri &
biên dịch / lúc chạy / không phát hiện \\
RQ4 & Nguyên mẫu di trú FFI & Máy chủ C++ $+$ đường nối C-ABI Rust &
kích thước đường nối, phạm vi \texttt{unsafe}, chi phí mỗi lần gọi \\
\hline
\end{tabular}
\end{table}

\section[Phương pháp theo từng câu hỏi]{Phương pháp theo từng câu hỏi nghiên cứu, tóm tắt}\label{sec:methodbrief}
RQ1 và RQ3 được trả lời bằng một quét tải vòng hở từ tải nhẹ đến vượt xa bão hoà. Vòng hở nghĩa
là bộ sinh giữ lịch gửi của nó thay vì chờ trả
lời~\cite{schroeder2006openclosed}. Bốn biến thể build được đan xen trong mỗi lần lặp. Một phép
đo thời gian thứ hai tách thao tác lõi khỏi bàn giao giữa luồng quanh nó. Vòng khứ hồi đầu-cuối
gói hai chi phí thuộc hai loại. Chỉ một trong hai có thể khác giữa các ngôn ngữ. Một phép đo
thứ ba lấy hàng đợi tranh chấp và tay cầm trả lời riêng. Vì vậy nguyên nhân quá tải có thể được
thử chứ không chỉ được lập luận.

RQ2 được trả lời bằng tiêm lỗi có kiểm soát~\cite{natella2016fault}. Mỗi lớp lỗi là một chương
trình nhỏ riêng ở cả hai ngôn ngữ. Phía C++ được xây và chạy dưới AddressSanitizer. Nó cũng
chạy dưới ThreadSanitizer và UndefinedBehaviorSanitizer. Phía Rust an toàn được biên dịch và,
nơi nó biên dịch, chạy dưới một hạn thời gian bị chặn. Không bộ phân tích tĩnh nào được chạy,
cố ý. Các cột C++ khi đó báo cáo công cụ mà nghiên cứu này đã chạy. Chúng không báo cáo trần
của những gì công cụ C++ có thể tìm. RQ4 được trả lời bằng một nguyên mẫu chạy được theo hướng
một cuộc di trú thực đi. Một máy chủ C++ gọi một lõi Rust không đổi qua ABI C. Chi phí mỗi lần
gọi được đo trên một thao tác thực chứ không suy từ một lần vượt trống. Lỗi bộ nhớ vượt ngôn
ngữ nằm ở đường nối đó~\cite{li2022crosslang}.

\section{Tính công bằng như một biến được kiểm soát}\label{sec:controlsbrief}
Một phép so sánh ngôn ngữ chỉ tốt bằng các biến nó giữ cố định. Đặc tả được cố định trước khi
bản nào tồn tại. Mô hình luồng, bộ cấp phát và số arena của nó, mức tối ưu, hàm băm, phương
pháp histogram và lịch tải được khớp giữa hai phía. Chúng không bao giờ được tinh chỉnh theo
từng phía. Nơi hai ngôn ngữ không thể viết giống nhau, khoảng cách được ghi lại chứ không được
hoà. Hai phía được chứng minh khớp bằng cách diff nhật ký khớp theo từng thao tác. Một hàng của
bảng đó không có đối chứng. Một người viết cả hai bản và không thành thạo như nhau ở cả hai.
Mục~\ref{sec:standing} coi đó là mối đe dọa nội tại chính của nghiên cứu.

\section{Xử lý thống kê}\label{sec:statsbrief}
Thiết kế thống kê được cố định trước. Họ các phép so sánh được thử được liệt kê trước chiến
dịch. Vì vậy một mức tải chọn sau khi nhìn một quét không thể vào họ. Mỗi phép so sánh là một
kiểm định hoán vị ghép cặp chính xác trên \NRepeats{} lần lặp đan xen, với khoảng tin cậy
bootstrap~\cite{efron1993bootstrap}, hiệu chỉnh Holm trên cả
họ~\cite{holm1979simple}. Các sàn mà kiểm định có thể đạt ($\MinAttainableP$ chưa hiệu chỉnh,
$\CorrectedFloor$ đã hiệu chỉnh) được nêu để một kết quả null phân biệt được với một giới hạn
toán. Hai đại lượng được giữ ngoài họ. Một là độ dao động từ lần gọi này sang lần gọi khác, mà
không câu hỏi nghiên cứu nào hỏi. Đại lượng kia là đuôi xa p99.99 gộp. Đơn vị của nó là thao
tác, không phải lần chạy. Nó khi đó không có khoảng giữa các lần lặp để thử.

\section{Tái lập và hiện vật nghiên cứu}\label{sec:repro}
Một bên thứ ba có thể xây lại kết quả mà không cần mã riêng của Nokia. Các bản được đưa chạy
hoàn toàn bởi đặc tả công khai đã cố định. Khung đo, danh mục lỗi và bộ chạy của nó, và nguyên
mẫu FFI được công bố cùng luận văn. Mỗi hình dữ liệu đo được một kịch bản xây lại từ nhật ký
thô. Mỗi số được trích trong kết quả cũng vậy. Kịch bản đó là
\texttt{bench/refresh\_thesis.sh}. Tuyên bố đó không phủ các sơ đồ vẽ tay, vốn giải thích một
khái niệm chứ không báo cáo một kết quả. Và chạy cùng các kịch bản trên chiến dịch của một máy
khác xây lại hiện vật theo cùng thủ tục, không với cùng giá trị.

Các nhật ký mà kịch bản đọc là một phần của gói. Mỗi chiến dịch có kịch bản riêng và tiền tố
macro riêng, nên một hình từ chiến dịch này không thể chạm một câu về chiến dịch kia. Thiết lập
máy chủ của mỗi lần chạy được ghi và lưu cạnh dữ liệu chiến dịch. Không gì trong các chương kết
quả được gõ từ những tệp đó bằng tay.
